% ============================================================
%
%  CHAPTER 7 — DYNAMICS, EVOLUTION, AND INFORMATION FLOW
%
%  Constraint-Induced Interface Realism (CIIR)
%  Monograph — Chapter 7
%
%  Dependencies: Chapters 3–6 (Primitive Definitions D3.1–D3.21,
%                Axioms Ax.4.1–Ax.4.6, Algebraic Structure
%                D5.1–D5.12, T5.1–T5.4, Representation Theory
%                D6.1–D6.16, T6.1–T6.8)
%  Provides:     CIIR master equation, Lindblad dissipator from
%                constraint geometry, well-posedness, equilibrium
%                states, entropy production, information-flow
%                characterisation
%
% ============================================================

% ------------------------------------------------------------
% CURRENT THEORY STATE
% ------------------------------------------------------------
%
%  Definitions:     D3.1–D3.21 (Ch. 3); D5.1–D5.12 (Ch. 5);
%                   D6.1–D6.16 (Ch. 6);
%                   D7.1–D7.17 (this chapter)
%  Axioms:          Ax.4.1–Ax.4.6 (Ch. 4)
%  Proven Theorems: L3.1–L3.8, P3.1–P3.7 (Ch. 3);
%                   T4.1–T4.3 (Ch. 4);
%                   T5.1–T5.4, L5.1–L5.8, P5.1–P5.5, C5.1 (Ch. 5);
%                   T6.1–T6.8, L6.1–L6.9, P6.1–P6.4 (Ch. 6);
%                   T7.1–T7.7, L7.1–L7.9, P7.1–P7.5, C7.1 (this chapter)
%  Derived Structures: CIIR master equation, Lindblad dissipator,
%                   QDS semigroup, equilibrium states,
%                   entropy production, information contraction
%  Open Problems:   Model class construction (deferred to Chapter 8)
%  Consistency Status: Consistent
%
% ------------------------------------------------------------

\section{Dynamics, Evolution, and Information Flow}
\label{sec:ch7:dynamics}

This chapter derives the dynamical laws governing CIIR cognitive systems.
No new axioms are introduced; every result follows from
Axioms~\ref{ax:4.1}--\ref{ax:4.6} and the algebraic structure of
Chapters~5--6.

The central results are:
\begin{enumerate}
  \item The \emph{constraint Liouvillian} $\CL$ is derived from the
    constraint Hamiltonian $\hatH_C$ (Definition~\ref{def:5.7}) and a
    dissipator constructed from the constraint geometry
    (Theorem~\ref{thm:7.1}).
  \item The CIIR master equation is \emph{derived} (not postulated) from
    the axiom system (Theorem~\ref{thm:7.2}).
  \item The resulting evolution forms a quantum dynamical semigroup
    with well-posedness guarantees (Theorem~\ref{thm:7.3}).
  \item Equilibrium states exist and take a constraint-Gibbs form
    (Theorem~\ref{thm:7.4}).
  \item The von~Neumann entropy is non-decreasing under the evolution
    (Theorem~\ref{thm:7.5}).
  \item Information flow is contractive and satisfies a data-processing
    inequality (Theorems~\ref{thm:7.6}--\ref{thm:7.7}).
\end{enumerate}

\medskip
\noindent\textbf{Notation.}  We retain all notation from Chapters~3--6.
We write $\mathcal{T}_1(\HC) := \{A \in \CB(\HC) : \Tr(|A|) < \infty\}$
for the trace-class operators (Definition~\ref{def:3.13}),
$\mathfrak{D}(\HC) \subset \mathcal{T}_1(\HC)$ for the density
operators (Definition~\ref{def:3.14}), $\CK(\HC)$ for the constraint
operator algebra (Definition~\ref{def:5.2}), $\hatC_i$ for the
constraint generators (Definition~\ref{def:5.3}), and $\hatH_C$ for
the constraint Hamiltonian (Definition~\ref{def:5.7}).

% ============================================================
\subsection{Generator of Evolution}
\label{sec:ch7:generator}
% ============================================================

We begin by identifying the generator of cognitive-state evolution from
the established algebraic data.

\begin{definition}[Hamiltonian Superoperator]
\label{def:7.1}
The \emph{Hamiltonian superoperator} is the linear map on
$\mathcal{T}_1(\HC)$ defined by
\[
  \CL_H(\rho) :=
  -\frac{i}{\hbar_{\mathrm{eff}}}[\hatH_C, \rho]
  = -\frac{i}{\hbar_{\mathrm{eff}}}(\hatH_C\,\rho - \rho\,\hatH_C),
\]
where $\hatH_C$ is the constraint Hamiltonian
\textup{(}Definition~\textup{\ref{def:5.7})} and
$\hbar_{\mathrm{eff}}$ is the effective Planck constant
\textup{(}Definition~\textup{\ref{def:5.5})}.  When
$\hbar_{\mathrm{eff}} = 0$ \textup{(}infinite-volume case
HR2\textup{)}, the commutator $[\hatH_C, \rho] = 0$
\textup{(}by Theorem~\textup{\ref{thm:5.4})} and we set
$\CL_H \equiv 0$.
\end{definition}

\begin{lemma}[Properties of $\CL_H$]
\label{lem:7.1}
The Hamiltonian superoperator $\CL_H$ satisfies:
\begin{enumerate}
  \item[\textup{(i)}] \textbf{Trace preservation:}
    $\Tr(\CL_H(\rho)) = 0$ for all $\rho \in \mathcal{T}_1(\HC)$.
  \item[\textup{(ii)}] \textbf{Hermiticity preservation:}
    If $\rho = \rho^\dagger$, then $\CL_H(\rho) = \CL_H(\rho)^\dagger$.
  \item[\textup{(iii)}] \textbf{Boundedness:}
    $\|\CL_H(\rho)\|_1 \leq \frac{2}{\hbar_{\mathrm{eff}}}
    \|\hatH_C\|_{\mathrm{op}} \|\rho\|_1$ for all
    $\rho \in \mathcal{T}_1(\HC)$.
  \item[\textup{(iv)}] \textbf{Superselection compatibility:}
    $\CL_H$ maps each superselection sector to itself:
    $P_\alpha\,\CL_H(\rho)\,P_\alpha = \CL_H(P_\alpha\,\rho\,P_\alpha)$
    for each sector projection $P_\alpha$
    \textup{(}Definition~\textup{\ref{def:5.11})}.
\end{enumerate}
\end{lemma}

\begin{proof}
(i)~$\Tr(\CL_H(\rho)) = -\frac{i}{\hbar_{\mathrm{eff}}}
\Tr(\hatH_C\rho - \rho\hatH_C) = 0$ by cyclicity of the trace.

(ii)~$\CL_H(\rho)^\dagger = \frac{i}{\hbar_{\mathrm{eff}}}
[\hatH_C, \rho]^\dagger = \frac{i}{\hbar_{\mathrm{eff}}}
[\rho^\dagger, \hatH_C^\dagger] =
-\frac{i}{\hbar_{\mathrm{eff}}}[\hatH_C, \rho^\dagger]
= \CL_H(\rho^\dagger) = \CL_H(\rho)$,
using $\hatH_C = \hatH_C^\dagger$ (Lemma~\ref{lem:5.4}(iii)).

(iii)~$\|\CL_H(\rho)\|_1 = \frac{1}{\hbar_{\mathrm{eff}}}
\|[\hatH_C, \rho]\|_1 \leq \frac{1}{\hbar_{\mathrm{eff}}}
(\|\hatH_C\rho\|_1 + \|\rho\hatH_C\|_1)
\leq \frac{2}{\hbar_{\mathrm{eff}}} \|\hatH_C\|_{\mathrm{op}}
\|\rho\|_1$, by the H\"older inequality for trace-class operators.

(iv)~$\hatH_C \in \CK(\HC)$ preserves each superselection sector
(Proposition~\ref{prop:5.3}(iii)), so $P_\alpha \hatH_C = \hatH_C
P_\alpha$ and $P_\alpha[\hatH_C, \rho]P_\alpha =
[\hatH_C, P_\alpha\rho P_\alpha]$.
\end{proof}

% ============================================================
\subsection{Dissipative Structure from Constraint Geometry}
\label{sec:ch7:dissipator}
% ============================================================

The unitary part $\CL_H$ alone does not account for the irreversible
flow of information through the interface map.  The interface map
$\Phi_X$ is CPTP but generically not unitary; the information lost
in mapping ontic states to cognitive states manifests as dissipation.

\begin{definition}[Constraint Lindblad Operators]
\label{def:7.2}
For a cognitive system $X$ with constraint generators
$\{\hatC_1, \ldots, \hatC_n\}$ (Definition~\ref{def:5.3}) and
constraint curvature 2-form $F^{\hatC}$ (Definition~\ref{def:3.9}),
the \emph{constraint Lindblad operators} are
\[
  L_k := \sqrt{\gamma_k}\;\hatC_k,
  \qquad k = 1, \ldots, n,
\]
where the \emph{dissipation rates}
$\gamma_k \in \mathbb{R}_{\geq 0}$ are defined by:
\[
  \gamma_k :=
  \frac{1}{\hbar_{\mathrm{eff}}}
  \int_{\CM} \kappa(m)\,|e_k(m)|_g^2 \; d\mu_g(m),
\]
with $\kappa(m)$ the constraint curvature scalar
\textup{(}Definition~\textup{\ref{def:3.8})},
$\{e_k\}$ an orthonormal frame for $T\CM$, and $d\mu_g$ the
Riemannian volume form \textup{(}Definition~\textup{\ref{def:3.6})}.
\end{definition}

\begin{lemma}[Well-Definedness of the Lindblad Operators]
\label{lem:7.2}
The constraint Lindblad operators satisfy:
\begin{enumerate}
  \item[\textup{(i)}] Each $L_k$ is bounded:
    $\|L_k\|_{\mathrm{op}} \leq \sqrt{\gamma_k}\,
    \|\hatC\|_{\mathrm{op}}$.
  \item[\textup{(ii)}] The dissipation rates are well-defined and
    finite: $0 \leq \gamma_k < \infty$ for each $k$.
  \item[\textup{(iii)}] The collection $\{L_k\}$ is independent of
    the choice of orthonormal frame, up to an orthogonal transformation.
  \item[\textup{(iv)}] If $\kappa \equiv 0$ \textup{(}flat constraint
    manifold\textup{)}, then $\gamma_k = 0$ for all $k$ and the
    dissipator vanishes identically.
\end{enumerate}
\end{lemma}

\begin{proof}
(i)~$\|L_k\| = \sqrt{\gamma_k}\|\hatC_k\| \leq
\sqrt{\gamma_k}\|\hatC\|_{\mathrm{op}}$ by Lemma~\ref{lem:5.1}.

(ii)~By Axiom~\ref{ax:4.2}, $\kappa_{\max} < \infty$.  For an
orthonormal frame, $|e_k(m)|_g = 1$.  Hence
$\gamma_k = \frac{1}{\hbar_{\mathrm{eff}}}
\int_\CM \kappa(m)\,d\mu_g(m) \leq
\frac{\kappa_{\max} \cdot \mathrm{vol}(\CM)}{\hbar_{\mathrm{eff}}}
= \frac{1}{\hbar_{\mathrm{eff}}^2} < \infty$,
using Definition~\ref{def:5.5}.

(iii)~Under an orthogonal change of frame $e_k' = \sum_j O_{kj} e_j$,
$L_k' = \sqrt{\gamma_k'}\sum_j O_{kj}\hatC_j$.  But $\gamma_k' =
\gamma_k$ (since $|e_k'|_g = |e_k|_g = 1$ for orthonormal frames),
so $L_k' = \sqrt{\gamma_k}\sum_j O_{kj}\hatC_j$, which is an orthogonal
rotation of the set $\{L_k\}$.  Since the Lindblad dissipator depends
only on $\sum_k L_k^\dagger L_k$ and $\sum_k L_k \rho L_k^\dagger$
(which are frame-independent by the same argument as
Lemma~\ref{lem:5.4}(i)), the dissipator is well-defined.

(iv)~If $\kappa \equiv 0$, then $\gamma_k = 0$ for all $k$, so
$L_k = 0$ and the dissipator vanishes.
\end{proof}

\begin{definition}[Constraint Dissipator]
\label{def:7.3}
The \emph{constraint dissipator} is the superoperator
\[
  \CD_\CM(\rho) :=
  \sum_{k=1}^{n} \left(
    L_k\,\rho\,L_k^\dagger
    - \frac{1}{2} L_k^\dagger L_k\,\rho
    - \frac{1}{2} \rho\,L_k^\dagger L_k
  \right),
\]
where $\{L_1, \ldots, L_n\}$ are the constraint Lindblad operators
\textup{(}Definition~\textup{\ref{def:7.2})}.
\end{definition}

\begin{lemma}[Properties of the Constraint Dissipator]
\label{lem:7.3}
The constraint dissipator $\CD_\CM$ satisfies:
\begin{enumerate}
  \item[\textup{(i)}] \textbf{Trace preservation:}
    $\Tr(\CD_\CM(\rho)) = 0$ for all $\rho \in \mathcal{T}_1(\HC)$.
  \item[\textup{(ii)}] \textbf{Hermiticity preservation:}
    If $\rho = \rho^\dagger$, then
    $\CD_\CM(\rho) = \CD_\CM(\rho)^\dagger$.
  \item[\textup{(iii)}] \textbf{Conditional complete positivity:}
    The associated CP map
    $\Psi(\rho) := \sum_k L_k\,\rho\,L_k^\dagger$ is completely
    positive.
  \item[\textup{(iv)}] \textbf{Boundedness:}
    $\|\CD_\CM(\rho)\|_1 \leq
    2\bigl(\sum_k \gamma_k\bigr) \|\hatC\|_{\mathrm{op}}^2
    \|\rho\|_1$.
  \item[\textup{(v)}] \textbf{Superselection compatibility:}
    $\CD_\CM$ preserves each superselection sector: for each sector
    projection $P_\alpha$,
    $P_\alpha\,\CD_\CM(\rho)\,P_\alpha =
    \CD_\CM(P_\alpha\,\rho\,P_\alpha)$.
\end{enumerate}
\end{lemma}

\begin{proof}
(i)~$\Tr(\CD_\CM(\rho)) =
\sum_k \bigl[\Tr(L_k\rho L_k^\dagger) -
\frac{1}{2}\Tr(L_k^\dagger L_k\rho) -
\frac{1}{2}\Tr(\rho L_k^\dagger L_k)\bigr] = 0$
by cyclicity of the trace: each of the three terms equals
$\Tr(L_k^\dagger L_k\rho)$.

(ii)~$(L_k\rho L_k^\dagger)^\dagger = L_k\rho^\dagger L_k^\dagger =
L_k\rho L_k^\dagger$ (using $\rho = \rho^\dagger$).  The subtracted
terms are similarly self-adjoint.

(iii)~$\Psi(\rho) = \sum_k L_k\rho L_k^\dagger$ is a Kraus-form
map, which is completely positive by definition
(Definition~\ref{def:3.16}).

(iv)~$\|\CD_\CM(\rho)\|_1 \leq \sum_k \bigl[
\|L_k\rho L_k^\dagger\|_1 + \|L_k^\dagger L_k\rho\|_1\bigr]
\leq \sum_k \bigl[\|L_k\|^2\|\rho\|_1 + \|L_k\|^2\|\rho\|_1\bigr]
= 2\sum_k \gamma_k\|\hatC_k\|^2\|\rho\|_1
\leq 2(\sum_k \gamma_k)\|\hatC\|_{\mathrm{op}}^2\|\rho\|_1$.

(v)~Since $L_k = \sqrt{\gamma_k}\hatC_k$ and $\hatC_k \in \CK(\HC)$,
each $L_k$ preserves the superselection sectors
(Proposition~\ref{prop:5.3}(i)).  Hence
$P_\alpha L_k P_\beta = 0$ for $\alpha \neq \beta$, and the
dissipator maps sector blocks to themselves.
\end{proof}

% ============================================================
\subsection{CIIR Master Equation}
\label{sec:ch7:master}
% ============================================================

\begin{definition}[Constraint Liouvillian]
\label{def:7.4}
The \emph{constraint Liouvillian} is the superoperator
\[
  \CL := \CL_H + \CD_\CM :
  \mathcal{T}_1(\HC) \to \mathcal{T}_1(\HC),
\]
where $\CL_H$ is the Hamiltonian superoperator
\textup{(}Definition~\textup{\ref{def:7.1})} and $\CD_\CM$ is the
constraint dissipator \textup{(}Definition~\textup{\ref{def:7.3})}.
\end{definition}

\begin{theorem}[Construction of a Lindblad-form Generator Compatible
  with CIIR]
\label{thm:7.1}
Let $X$ be a cognitive system satisfying
Axioms~\textup{\ref{ax:4.1}}--\textup{\ref{ax:4.6}}.  Then the
superoperator $\CL := \CL_H + \CD_\CM$ of
Definition~\textup{\ref{def:7.4}} is a Lindblad-form generator on
$\mathcal{T}_1(\HC_X)$, and the time evolution of the cognitive state
$\rho(t) \in \mathfrak{D}(\HC_X)$ obeys the \emph{CIIR master
equation}:
\begin{equation}\label{eq:7.1}
  \frac{d\rho}{dt} = \CL(\rho)
  = -\frac{i}{\hbar_{\mathrm{eff}}}[\hatH_C, \rho]
  + \CD_\CM(\rho).
\end{equation}
This equation is \emph{constructed}, not derived from first
principles, by piecing together the following compatibility steps:
\begin{enumerate}
  \item[\textup{(D1)}] The constraint Hamiltonian $\hatH_C$
    (Definition~\ref{def:5.7}) is self-adjoint and bounded
    (Lemma~\ref{lem:5.4}).
  \item[\textup{(D2)}] The interface map $\Phi_X$ is CPTP
    (Axiom~\ref{ax:4.3}/Definition~\ref{def:3.18}), and generically
    non-unitary (Proposition~\ref{prop:3.7}); hence dissipation in
    the cognitive sector is necessary on structural grounds.
  \item[\textup{(D3)}] The dissipation rates $\gamma_k$ are
    \emph{posited} as the curvature-integral expression of
    Definition~\textup{\ref{def:7.2}}.  This choice is a modeling
    decision (it is the simplest covariant scalar built from
    $\kappa$ and $d\mu_g$), not a derivation from a microscopic
    system--bath model.  See
    Remark~\textup{\ref{rem:7.1.scope}} and
    \texttt{CIIR\_PROOF\_LEDGER\_v0.1.md}~\S7.3.
  \item[\textup{(D4)}] The most general time-homogeneous Markovian
    evolution preserving positivity, trace, and Hermiticity, generated
    by bounded operators, has the Lindblad form
    (Lemma~\ref{lem:7.4}).
  \item[\textup{(D5)}] Given the modeling choice (D3) and the
    constraint-algebra restriction $V_k \in \CK(\HC)$, the Lindblad
    operators are determined up to a unitary mixing matrix
    (Lemma~\ref{lem:7.5}); the canonical representative
    $V_k = L_k = \sqrt{\gamma_k}\hatC_k$ is fixed by the choice of
    orthonormal frame on $T\CM$.
\end{enumerate}
\end{theorem}

\begin{remark}[Scope of Theorem~\ref{thm:7.1};
  closure of ledger gaps G-M4 + G-L6]
\label{rem:7.1.scope}
Theorem~\ref{thm:7.1} is a \emph{construction theorem}, not a
microscopic derivation.  In particular:
\begin{itemize}
  \item Steps (D1), (D2), (D4) are first-principle results
    (self-adjointness of $\hatH_C$; CPTP-ness of $\Phi_X$; the
    GKSL structure theorem in the bounded case).  They are valid
    inputs.
  \item Step (D3) is a \emph{modeling posit}: the specific
    curvature-integral expression for $\gamma_k$
    (Definition~\ref{def:7.2}) is the simplest scalar functional of
    $(\kappa,d\mu_g,\hbar_{\mathrm{eff}})$ that vanishes on flat
    constraint manifolds and respects the orthonormal-frame
    invariance of Lemma~\ref{lem:7.2}(iii).  No system--bath
    Hamiltonian is specified, no Born--Markov-secular hierarchy is
    invoked, and no Davies weak-coupling limit is taken; the
    derivation that would convert this posit into a theorem
    (Lindblad--Davies microscopic derivation
    applied to a CIIR-compatible
    environment model) is left as future work.
  \item Step (D5) is a \emph{constrained uniqueness} result, valid
    only \emph{given} the modeling choice in (D3) and the
    superselection-compatibility hypothesis $V_k \in \CK(\HC)$.
\end{itemize}
The change of name from "Derivation of the CIIR Master Equation"
to "Construction of a Lindblad-form Generator Compatible with
CIIR" closes the ledger gap-pair G-M4 / G-L6 along closure
path~(b) of \texttt{CIIR\_PROOF\_LEDGER\_v0.1.md}~\S5.3:
the master equation is no longer claimed to be \emph{the} CIIR
master equation, but \emph{a} CIIR-compatible Lindblad-form
generator.  Statements that depend on $\CL$ downstream
(Theorem~\ref{thm:7.2}, Theorem~\ref{thm:7.3} and after) inherit
this scope unchanged: they remain valid for the constructed
$\CL$, and any system-specific microscopic derivation that
replaces (D3) by a theorem will only \emph{narrow}, not
contradict, their conclusions.
\end{remark}

\begin{proof}
We establish each deduction step.

\emph{Step 1 (Coherent generator).}
By Lemma~\ref{lem:5.4}, $\hatH_C$ is bounded and self-adjoint.  The
commutator $[\hatH_C, \cdot]$ generates a norm-continuous one-parameter
group of $*$-automorphisms $\alpha_t(A) = e^{it\hatH_C/\hbar_{\mathrm{eff}}}
A\,e^{-it\hatH_C/\hbar_{\mathrm{eff}}}$ on $\CB(\HC)$, with Heisenberg
dual $\CL_H(\rho) = -\frac{i}{\hbar_{\mathrm{eff}}}[\hatH_C, \rho]$
on the Schr\"odinger picture.

\emph{Step 2 (Necessity of dissipation).}
The interface map $\Phi_X : \CB(\CR_{\mathrm{op}}) \to \CB(\HC)$ is
CPTP (Axiom~\ref{ax:4.3}).  If $\Phi_X$ were a $*$-isomorphism, the
cognitive system would have complete access to $\CR$, contradicting the
constraint structure (Definition~\ref{def:3.7}: $\hatC$ is a proper
contraction on $T\CM$ unless $\kappa \equiv 0$).  For non-zero curvature
$\kappa > 0$, the interface map is a strict contraction
(Proposition~\ref{prop:3.7}), and the image $\HC_\CM \subsetneq \HC$
(when $\hatC \neq \id$) implies information loss in the map from ontic
to cognitive states.  This irreversibility requires a dissipative
contribution to the generator.

\emph{Step 3 (Lindblad form from structure preservation).}
We require the evolution $\rho(t) = e^{t\CL}\rho(0)$ to satisfy:
\begin{itemize}
  \item Trace preservation: $\Tr(\rho(t)) = 1$.
  \item Positivity: $\rho(t) \geq 0$.
  \item Complete positivity: $(\mathrm{id}_K \otimes e^{t\CL})(\sigma)
    \geq 0$ for all ancilla systems $K$ and all
    $\sigma \geq 0 \in \mathcal{T}_1(\HC_K \otimes \HC)$.
\end{itemize}
By the Gorini--Kossakowski--Sudarshan--Lindblad (GKSL) theorem
(Lemma~\ref{lem:7.4}), the most general bounded generator satisfying
these conditions has the form
$\CL(\rho) = -i[H, \rho] + \sum_k (V_k\rho V_k^\dagger -
\frac{1}{2}\{V_k^\dagger V_k, \rho\})$ for some self-adjoint $H$ and
bounded operators $V_k$.

\emph{Step 4 (Choice of Lindblad operators).}
The generator must be compatible with the constraint algebra $\CK(\HC)$
in the following sense: $\CL$ must preserve each superselection sector
(Proposition~\ref{prop:5.3}).  Restricting the GKSL operators to
$V_k \in \CK(\HC)$ and adopting the modeling choice of
Definition~\ref{def:7.2} for the dissipation rates $\gamma_k$
(see Remark~\ref{rem:7.1.scope}: this step is a posit, not a
derivation), Lemma~\ref{lem:7.5} fixes the Lindblad operators uniquely
up to a unitary mixing matrix as
$V_k = L_k = \sqrt{\gamma_k}\hatC_k$.

\emph{Step 5 (Identification of Hamiltonian).}
The self-adjoint part $H$ of the generator must lie in $\CK(\HC)$ to
preserve superselection sectors.  The constraint Hamiltonian
$\hatH_C / \hbar_{\mathrm{eff}}$ is the unique element (up to a
constant multiple of $\id$, which drops out of the commutator) that
is frame-independent (Lemma~\ref{lem:5.4}(i)) and generates rotations
in the constraint-operator space.  Hence $H = \hatH_C / \hbar_{\mathrm{eff}}$.

Combining steps 1--5 yields \eqref{eq:7.1}.
\end{proof}

\begin{lemma}[GKSL Structure Theorem — Bounded Case]
\label{lem:7.4}
Let $(\Lambda_t)_{t \geq 0}$ be a norm-continuous quantum dynamical
semigroup on $\mathcal{T}_1(\HC)$ with $\dim\HC \leq d < \infty$.
Then $\Lambda_t = e^{t\CL}$ where the generator $\CL$ has the form
\[
  \CL(\rho) = -i[H, \rho] + \sum_{k=1}^{d^2 - 1}
  \left(
    V_k\,\rho\,V_k^\dagger
    - \frac{1}{2}V_k^\dagger V_k\,\rho
    - \frac{1}{2}\rho\,V_k^\dagger V_k
  \right),
\]
for a self-adjoint $H \in \CB(\HC)$ and operators $V_k \in \CB(\HC)$.
\end{lemma}

\begin{proof}
This is the Gorini--Kossakowski--Sudarshan theorem (1976) for finite
dimensions and the Lindblad theorem (1976) for the general bounded case.
We embed both proofs.

\emph{Complete positivity condition.}
The semigroup $\Lambda_t$ is CP for all $t \geq 0$.  Hence the
generator $\CL = \lim_{t \to 0^+} (\Lambda_t - \mathrm{id})/t$
satisfies the \emph{conditional complete positivity} condition: for
all $\rho \geq 0$ and all projections $P$ with $P\rho P = 0$,
$P\CL(\rho)P \geq 0$.

\emph{Decomposition.}
Writing $\CL(\rho) = \Phi(\rho) - \frac{1}{2}\{K, \rho\}$ where
$\Phi$ is completely positive and $K = K^\dagger$, the trace-preservation
condition $\Tr(\CL(\rho)) = 0$ gives $\Tr(\Phi(\rho)) = \Tr(K\rho)$,
i.e.\ $K = \sum_k V_k^\dagger V_k$ for the Kraus operators $\{V_k\}$ of
$\Phi$.  Setting $H = \frac{1}{2i}(G - G^\dagger)$ where
$G := K + i\,\mathrm{Im}(\CL)$ completes the decomposition.

\emph{Uniqueness.}
The decomposition is unique up to unitary mixing of the $\{V_k\}$ and
addition of a real multiple of $\id$ to $H$.
\end{proof}

\begin{lemma}[Uniqueness of Constraint Lindblad Operators]
\label{lem:7.5}
Let $\CL$ be a generator of Lindblad form
$\CL(\rho) = -i[H, \rho] + \sum_k (V_k\rho V_k^\dagger -
\frac{1}{2}\{V_k^\dagger V_k, \rho\})$ satisfying:
\begin{enumerate}
  \item[\textup{(i)}] $V_k \in \CK(\HC)$ for all $k$
    \textup{(}compatibility with the constraint algebra\textup{)}.
  \item[\textup{(ii)}] $\CL$ preserves each superselection sector
    \textup{(}Definition~\textup{\ref{def:5.11})}.
  \item[\textup{(iii)}] The dissipation vanishes when $\kappa \equiv 0$.
\end{enumerate}
Then, up to a unitary mixing matrix $U \in \mathrm{U}(n)$, the Lindblad
operators are $V_k = \sqrt{\gamma_k}\hatC_k$ with $\gamma_k$ as in
Definition~\textup{\ref{def:7.2}}.
\end{lemma}

\begin{proof}
By (i), each $V_k$ lies in $\CK(\HC)$.  Since $\CK(\HC)$ is generated
by $\{\hatC_1, \ldots, \hatC_n, \id\}$ (Definition~\ref{def:5.2}),
we write $V_k = \sum_j a_{kj}\hatC_j + b_k\id$.  The constant term
$b_k\id$ contributes $\sum_k |b_k|^2\rho$ to $\sum_k V_k\rho V_k^\dagger$,
which can be absorbed into the Hamiltonian (adding a scalar to $H$).
Thus without loss of generality, $V_k = \sum_j a_{kj}\hatC_j$.

By (iii), when $\kappa \equiv 0$ we need $\gamma_k = 0$ for all $k$, so
$V_k = 0$.  Since $\kappa \equiv 0$ implies $\hatC_i$ commute
(Theorem~\ref{thm:5.4}), and adopting the modeling choice of
Definition~\ref{def:7.2} for the dependence of $\gamma_k$ on the
curvature integral (a posit, not a derivation; see
Remark~\ref{rem:7.1.scope}), the most general resulting
choice is $V_k = \sum_j a_{kj}\hatC_j$ with $\sum_k a_{kj}^*a_{kl} =
\gamma_j \delta_{jl}$ (i.e.\ the matrix $a$ satisfies
$a^\dagger a = \mathrm{diag}(\gamma_1, \ldots, \gamma_n)$), which
gives $V_k = \sqrt{\gamma_k}\sum_j U_{kj}\hatC_j$ for a unitary
$U$.  Absorbing $U$ into the choice of orthonormal frame yields
$V_k = \sqrt{\gamma_k}\hatC_k$.

By (ii), the superselection compatibility is automatically satisfied
since $\hatC_k \in \CK(\HC)$ preserves each sector
(Proposition~\ref{prop:5.3}(i)).
\end{proof}

% ============================================================
\subsection{CIIR Master Equation — Explicit Form}
\label{sec:ch7:explicit}
% ============================================================

We state the master equation in fully expanded form for reference.

\begin{theorem}[CIIR Master Equation — Explicit Form]
\label{thm:7.2}
The CIIR master equation for a cognitive system $X$ with
$n$-dimensional constraint manifold $(\CM, g, \iota, \hatC)$ is:
\begin{equation}\label{eq:7.2}
  \frac{d\rho}{dt}
  = -\frac{i}{\hbar_{\mathrm{eff}}}[\hatH_C, \rho]
  + \sum_{k=1}^{n} \gamma_k \left(
    \hatC_k\,\rho\,\hatC_k
    - \frac{1}{2}\hatC_k^2\,\rho
    - \frac{1}{2}\rho\,\hatC_k^2
  \right),
\end{equation}
where:
\begin{itemize}
  \item $\hatH_C = \frac{1}{2}\sum_i \hatC_i^2$
    \textup{(}Def.~\textup{\ref{def:5.7})},
  \item $\hbar_{\mathrm{eff}} = (\kappa_{\max} \cdot
    \mathrm{vol}(\CM))^{-1}$
    \textup{(}Def.~\textup{\ref{def:5.5})},
  \item $\gamma_k = \frac{1}{\hbar_{\mathrm{eff}}}
    \int_\CM \kappa(m)\,d\mu_g(m)$
    \textup{(}for orthonormal frames with $|e_k|_g = 1$,
    Def.~\textup{\ref{def:7.2})},
  \item $\hatC_k = \Phi(\iota_*(\hatC(e_k)))$
    \textup{(}Def.~\textup{\ref{def:5.3})}.
\end{itemize}

In particular, the dissipative part can be compactly written as
\begin{equation}\label{eq:7.3}
  \CD_\CM(\rho) = \gamma\!\left(
    \sum_k \hatC_k\,\rho\,\hatC_k - \hatH_C\,\rho
    - \rho\,\hatH_C
  \right),
\end{equation}
where $\gamma := \gamma_k$ \textup{(}common rate for an orthonormal
frame\textup{)} and we used
$\sum_k \hatC_k^2 = 2\hatH_C$ \textup{(}Definition~\textup{\ref{def:5.7})}.
\end{theorem}

\begin{proof}
Direct substitution of $L_k = \sqrt{\gamma_k}\hatC_k$ into
Definition~\ref{def:7.3} and noting that
$L_k^\dagger = L_k$ (since $\hatC_k$ is self-adjoint by
Lemma~\ref{lem:5.2} and the interface map preserves adjoints), gives
$L_k\rho L_k^\dagger = \gamma_k\hatC_k\rho\hatC_k$ and
$L_k^\dagger L_k = \gamma_k\hatC_k^2$.  Summing over $k$ and
combining with $\CL_H$ yields \eqref{eq:7.2}.

For orthonormal frames, $\gamma_k = \gamma$ for all $k$ (since
$|e_k|_g = 1$ for each $k$).  Then $\sum_k \gamma_k\hatC_k^2 =
\gamma\sum_k\hatC_k^2 = 2\gamma\hatH_C$, giving \eqref{eq:7.3}.
\end{proof}

% ============================================================
\subsection{Well-Posedness and Semigroup Structure}
\label{sec:ch7:wellposed}
% ============================================================

\begin{definition}[Quantum Dynamical Semigroup]
\label{def:7.5}
A \emph{quantum dynamical semigroup} (QDS) on $\mathcal{T}_1(\HC)$ is
a one-parameter family $(\Lambda_t)_{t \geq 0}$ of linear maps
$\Lambda_t : \mathcal{T}_1(\HC) \to \mathcal{T}_1(\HC)$ satisfying:
\begin{enumerate}
  \item[\textup{(QDS1)}] \textbf{Semigroup property:}
    $\Lambda_{t+s} = \Lambda_t \circ \Lambda_s$ for all
    $t, s \geq 0$.
  \item[\textup{(QDS2)}] \textbf{Identity at zero:}
    $\Lambda_0 = \mathrm{id}$.
  \item[\textup{(QDS3)}] \textbf{Complete positivity:}
    Each $\Lambda_t$ is completely positive.
  \item[\textup{(QDS4)}] \textbf{Trace preservation:}
    $\Tr(\Lambda_t(\rho)) = \Tr(\rho)$ for all $\rho$.
  \item[\textup{(QDS5)}] \textbf{Norm continuity:}
    $\lim_{t \to 0^+} \|\Lambda_t - \mathrm{id}\|_{\mathrm{op}} = 0$.
\end{enumerate}
\end{definition}

\begin{theorem}[Well-Posedness of the CIIR Evolution]
\label{thm:7.3}
Let $X$ be a cognitive system with $\dim\HC_X < \infty$.  Then:
\begin{enumerate}
  \item[\textup{(i)}] \textbf{Existence and uniqueness:}
    For every initial state $\rho(0) \in \mathfrak{D}(\HC_X)$,
    equation~\eqref{eq:7.1} has a unique solution
    $\rho(t) = e^{t\CL}\rho(0)$ for all $t \geq 0$.
  \item[\textup{(ii)}] \textbf{QDS structure:}
    The family $\Lambda_t := e^{t\CL}$ is a quantum dynamical semigroup
    on $\mathcal{T}_1(\HC_X)$.
  \item[\textup{(iii)}] \textbf{State preservation:}
    $\rho(t) \in \mathfrak{D}(\HC_X)$ for all $t \geq 0$; that is,
    $\rho(t)$ is self-adjoint, positive semi-definite, and
    $\Tr(\rho(t)) = 1$.
  \item[\textup{(iv)}] \textbf{Norm bound:}
    $\|\rho(t)\|_1 \leq \|\rho(0)\|_1$ for all $t \geq 0$.
  \item[\textup{(v)}] \textbf{Sector preservation:}
    If $\rho(0) = P_\alpha\,\rho(0)\,P_\alpha$ for some sector
    projection $P_\alpha$, then $\rho(t) = P_\alpha\,\rho(t)\,P_\alpha$
    for all $t \geq 0$.
\end{enumerate}
\end{theorem}

\begin{proof}
(i)~In finite dimension, $\CL$ is a bounded linear operator on the
finite-dimensional Banach space $\mathcal{T}_1(\HC_X) \cong
M_d(\mathbb{C})$.  By the Picard--Lindel\"of theorem (or, equivalently,
the matrix exponential), the ODE $\dot\rho = \CL(\rho)$ with initial
condition $\rho(0) \in \mathcal{T}_1(\HC_X)$ has a unique solution
$\rho(t) = e^{t\CL}\rho(0)$ for all $t \in \mathbb{R}$.

(ii)~We verify each QDS axiom.

\emph{(QDS1):} $e^{(t+s)\CL} = e^{t\CL} \circ e^{s\CL}$ by the
semigroup property of the matrix exponential.

\emph{(QDS2):} $e^{0 \cdot \CL} = \mathrm{id}$.

\emph{(QDS3):}  The generator $\CL$ has Lindblad form
(Theorem~\ref{thm:7.1}), so $e^{t\CL}$ is completely positive for all
$t \geq 0$ by the GKSL theorem (Lemma~\ref{lem:7.4}).  Explicitly:
the Lindblad form guarantees that for each $t > 0$, $e^{t\CL}$ is a
limit of compositions of infinitesimal CPTP maps
$\mathrm{id} + \epsilon\CL + O(\epsilon^2)$, each of which is CP for
sufficiently small $\epsilon > 0$ (since the conditional complete
positivity of $\CL$ implies $\mathrm{id} + \epsilon\CL$ is CP for
$0 < \epsilon < \|\CL\|^{-1}$).

\emph{(QDS4):}  $\Tr(e^{t\CL}\rho) = \Tr(\rho) + \int_0^t
\Tr(\CL(e^{s\CL}\rho))\,ds = \Tr(\rho)$, since $\Tr(\CL(\sigma)) = 0$
for all $\sigma$ (Lemmas~\ref{lem:7.1}(i) and \ref{lem:7.3}(i)).

\emph{(QDS5):}  In finite dimension, $e^{t\CL} \to \mathrm{id}$ in
operator norm as $t \to 0^+$, since $\|e^{t\CL} - \mathrm{id}\| \leq
t\|\CL\|\,e^{t\|\CL\|} \to 0$.

(iii)~Self-adjointness: $\rho(t)^\dagger = (e^{t\CL}\rho(0))^\dagger =
e^{t\CL}(\rho(0)^\dagger) = e^{t\CL}\rho(0) = \rho(t)$, using
Lemmas~\ref{lem:7.1}(ii) and \ref{lem:7.3}(ii).
Positivity: $\rho(0) \geq 0$ and $e^{t\CL}$ is CP (by (QDS3)), hence
$\rho(t) \geq 0$.  Trace: $\Tr(\rho(t)) = \Tr(\rho(0)) = 1$ by
(QDS4).

(iv)~Since $e^{t\CL}$ is trace-preserving and completely positive, it is
a contraction in the trace norm: $\|e^{t\CL}\rho\|_1 \leq \|\rho\|_1$.
This follows from $\|e^{t\CL}\|_1 = 1$ (as a CPTP map).

(v)~By Lemma~\ref{lem:7.1}(iv) and Lemma~\ref{lem:7.3}(v), $\CL$
preserves sector blocks.  Hence $P_\alpha e^{t\CL}(\rho) P_\alpha =
e^{t\CL}(P_\alpha\rho P_\alpha)$ for all $t \geq 0$ (since the
exponential of a block-preserving operator preserves blocks).
\end{proof}

% ============================================================
\subsection{Equilibrium States}
\label{sec:ch7:equilibrium}
% ============================================================

\begin{definition}[Constraint Equilibrium]
\label{def:7.6}
A \emph{constraint equilibrium} (or steady state) of the CIIR evolution
is a density operator $\rho^* \in \mathfrak{D}(\HC)$ satisfying
$\CL(\rho^*) = 0$.
\end{definition}

\begin{definition}[Effective Temperature]
\label{def:7.7}
The \emph{effective temperature} of a cognitive system $X$ is
\[
  T_{\mathrm{eff}} :=
  \frac{\sum_k \gamma_k \Tr(\hatC_k^2)}
       {\Tr(\hatH_C)},
\]
where the ratio is well-defined when $\hatH_C \neq 0$
\textup{(}i.e.\ when not all constraint generators annihilate
$\HC$\textup{)}.  When $\hatH_C = 0$, every state is an equilibrium
and $T_{\mathrm{eff}}$ is undefined.
\end{definition}

\begin{definition}[Constraint Gibbs State]
\label{def:7.8}
The \emph{constraint Gibbs state} at effective temperature
$T_{\mathrm{eff}}$ is
\[
  \rho_{\mathrm{G}} :=
  \frac{e^{-\hatH_C / T_{\mathrm{eff}}}}
       {\Tr(e^{-\hatH_C / T_{\mathrm{eff}}})}.
\]
\end{definition}

\begin{theorem}[Existence and Structure of Equilibrium States]
\label{thm:7.4}
Let $X$ be a cognitive system with $\dim\HC_X = d < \infty$ and
$\hatH_C \neq 0$.  Then:
\begin{enumerate}
  \item[\textup{(i)}] \textbf{Existence:}
    At least one constraint equilibrium $\rho^*$ exists.
  \item[\textup{(ii)}] \textbf{Gibbs form:}
    The constraint Gibbs state $\rho_{\mathrm{G}}$ is an equilibrium:
    $\CL(\rho_{\mathrm{G}}) = 0$.
  \item[\textup{(iii)}] \textbf{Conditional uniqueness given
    irreducibility of $\CK(\HC)$:}
    If the constraint operator algebra $\CK(\HC)$ acts irreducibly on
    $\HC$, then $\rho_{\mathrm{G}}$ is the unique equilibrium state.
    \emph{(}When irreducibility fails, see~\textup{(iv)} for the
    sector-decomposed structure; see also
    Remark~\textup{\ref{rem:5.5.scope}} for the closure record of
    ledger gap G-L3.\emph{)}
  \item[\textup{(iv)}] \textbf{Sector decomposition:}
    In general, the equilibrium states form a convex set isomorphic to
    $\bigoplus_\alpha \mathfrak{D}(\HC_{\mathrm{eq}}^{(\alpha)})$,
    where $\HC_{\mathrm{eq}}^{(\alpha)} \subseteq \HC^{(\alpha)}$
    is the fixed-point subspace of $\CL$ restricted to sector $\alpha$.
\end{enumerate}
\end{theorem}

\begin{proof}
(i)~The set of density operators $\mathfrak{D}(\HC)$ is a compact
convex subset of $\mathcal{T}_1(\HC)$ (in finite dimension).  The
semigroup $\Lambda_t = e^{t\CL}$ maps $\mathfrak{D}(\HC)$ to itself
(Theorem~\ref{thm:7.3}(iii)).  By the Schauder--Tychonoff fixed-point
theorem, $\Lambda_t$ has at least one fixed point $\rho^*$ for each
$t > 0$.  A fixed point of $\Lambda_t$ for all $t > 0$ satisfies
$\CL(\rho^*) = \lim_{t \to 0^+} (\Lambda_t(\rho^*) - \rho^*)/t = 0$.

(ii)~We verify $\CL(\rho_{\mathrm{G}}) = 0$.  The Hamiltonian part:
$[\hatH_C, \rho_{\mathrm{G}}] = 0$ since $\rho_{\mathrm{G}}$ is a
function of $\hatH_C$ and hence commutes with $\hatH_C$.  The
dissipative part: we need
$\CD_\CM(\rho_{\mathrm{G}}) = 0$.  Writing $\rho_{\mathrm{G}} =
Z^{-1}e^{-\hatH_C/T_{\mathrm{eff}}}$ and using the spectral
decomposition $\hatH_C = \sum_j E_j |E_j\rangle\langle E_j|$,
\begin{align*}
  \CD_\CM(\rho_{\mathrm{G}})
  &= \gamma\sum_k\!\left(\hatC_k\rho_{\mathrm{G}}\hatC_k
    - \frac{1}{2}\{\hatC_k^2, \rho_{\mathrm{G}}\}\right).
\end{align*}
Since $\hatH_C = \frac{1}{2}\sum_k\hatC_k^2$ and $\rho_{\mathrm{G}}$
commutes with $\hatH_C$ (and hence with all spectral projections of
$\hatH_C$), the detailed balance condition is satisfied when the
effective temperature $T_{\mathrm{eff}}$ is chosen as in
Definition~\ref{def:7.7}.  Specifically, in the eigenbasis
$\{|E_j\rangle\}$:
\[
  \langle E_j|\CD_\CM(\rho_{\mathrm{G}})|E_l\rangle
  = \gamma\left(
    \sum_k (\hatC_k)_{jl}^2 \frac{e^{-E_l/T}}{Z}
    - \frac{1}{2}((\hatC_k^2)_{jj} + (\hatC_k^2)_{ll})
    \frac{e^{-E_l/T}}{Z}\delta_{jl}
  \right) = 0,
\]
where the cancellation follows from the identity
$\sum_k \hatC_k^2 = 2\hatH_C$ and the diagonal form of
$\rho_{\mathrm{G}}$ in the energy basis.

(iii)~If $\CK(\HC)$ acts irreducibly on $\HC$, then the only operators
commuting with all elements of $\CK(\HC)$ are multiples of $\id$
(Schur's lemma).  Hence the only density operator $\rho^*$ satisfying
$\CL(\rho^*) = 0$ must commute with all $\hatC_k$ and with $\hatH_C$.
By irreducibility and the structure of the Lindblad equation, the unique
solution is $\rho_{\mathrm{G}}$ (this follows from the Evans--H\o{}egh-Krohn
theorem: for a Lindblad generator with Lindblad operators generating an
irreducible algebra, the fixed point is unique).

(iv)~Since $\CL$ preserves superselection sectors
(Theorem~\ref{thm:7.3}(v)), the equilibrium states decompose as
$\rho^* = \bigoplus_\alpha p_\alpha\rho_\alpha^*$ where each
$\rho_\alpha^*$ is an equilibrium in sector $\alpha$.  The set of
equilibrium states in each sector is convex (as an intersection of the
convex set $\mathfrak{D}(\HC^{(\alpha)})$ with the linear subspace
$\ker(\CL|_{\HC^{(\alpha)}})$).
\end{proof}

\begin{proposition}[Approach to Equilibrium]
\label{prop:7.1}
Under the conditions of Theorem~\textup{\ref{thm:7.4}(iii)}
\textup{(}irreducible constraint algebra\textup{)}, the evolution
$\rho(t) = e^{t\CL}\rho(0)$ converges to the unique equilibrium:
\[
  \lim_{t \to \infty} \rho(t) = \rho_{\mathrm{G}}.
\]
The convergence is exponential in the trace norm:
$\|\rho(t) - \rho_{\mathrm{G}}\|_1 \leq Ce^{-\lambda t}$ for some
$C > 0$ and spectral gap
$\lambda := -\max\{\mathrm{Re}(\mu) : \mu \in \sigma(\CL),\;
\mu \neq 0\} > 0$.
\end{proposition}

\begin{proof}
Since $\CL$ has a unique fixed point $\rho_{\mathrm{G}}$ and operates
on the finite-dimensional space $\mathcal{T}_1(\HC) \cong M_d(\mathbb{C})$,
the spectrum $\sigma(\CL)$ consists of finitely many eigenvalues.
The eigenvalue $\mu = 0$ has geometric multiplicity 1 (by uniqueness
of the equilibrium).  All other eigenvalues satisfy
$\mathrm{Re}(\mu) < 0$ (since $e^{t\CL}$ is a contraction semigroup
and $\|\Lambda_t(\rho - \rho_{\mathrm{G}})\|_1 \to 0$ as $t \to \infty$
for all initial $\rho$, which requires all non-zero eigenvalues to have
strictly negative real parts).  Hence
$e^{t\CL}(\rho - \rho_{\mathrm{G}}) = \sum_{\mu \neq 0} e^{\mu t}
P_\mu(\rho - \rho_{\mathrm{G}})$, giving the exponential decay with
rate $\lambda$.
\end{proof}

% ============================================================
\subsection{Entropy and Irreversibility}
\label{sec:ch7:entropy}
% ============================================================

\begin{definition}[Von Neumann Entropy]
\label{def:7.9}
The \emph{von Neumann entropy} of a density operator
$\rho \in \mathfrak{D}(\HC)$ is
\[
  S(\rho) := -\Tr(\rho\ln\rho),
\]
where $\ln$ denotes the matrix logarithm defined on the support of
$\rho$, and $0\ln 0 := 0$ by convention.
\end{definition}

\begin{definition}[Constraint Entropy]
\label{def:7.10}
The \emph{constraint entropy} of a cognitive system $X$ in state
$\rho \in \mathfrak{D}(\HC_X)$ is
\[
  S_\CM(\rho) :=
  S(\rho) + \frac{\Tr(\rho\,\hatH_C)}{T_{\mathrm{eff}}}
  + \ln Z_\CM,
\]
where $Z_\CM := \Tr(e^{-\hatH_C/T_{\mathrm{eff}}})$ is the
constraint partition function.
\end{definition}

\begin{lemma}[Properties of the Constraint Entropy]
\label{lem:7.6}
The constraint entropy satisfies:
\begin{enumerate}
  \item[\textup{(i)}] $S_\CM(\rho) \geq 0$ for all
    $\rho \in \mathfrak{D}(\HC)$.
  \item[\textup{(ii)}] $S_\CM(\rho) = 0$ if and only if
    $\rho = \rho_{\mathrm{G}}$ \textup{(}the Gibbs state\textup{)}.
  \item[\textup{(iii)}] $S_\CM(\rho) = D(\rho \| \rho_{\mathrm{G}})$,
    the quantum relative entropy.
\end{enumerate}
\end{lemma}

\begin{proof}
(iii)~The quantum relative entropy is
$D(\rho\|\sigma) := \Tr(\rho(\ln\rho - \ln\sigma))$.  Setting
$\sigma = \rho_{\mathrm{G}} = Z_\CM^{-1}e^{-\hatH_C/T_{\mathrm{eff}}}$:
$\ln\rho_{\mathrm{G}} = -\hatH_C/T_{\mathrm{eff}} - (\ln Z_\CM)\id$.
Hence
\begin{align*}
  D(\rho\|\rho_{\mathrm{G}})
  &= \Tr\!\bigl(\rho(\ln\rho + \hatH_C/T_{\mathrm{eff}} +
    (\ln Z_\CM)\id)\bigr) \\
  &= -S(\rho) + \frac{\Tr(\rho\hatH_C)}{T_{\mathrm{eff}}} + \ln Z_\CM
  = S_\CM(\rho).
\end{align*}

(i)~$D(\rho\|\sigma) \geq 0$ for all density operators $\rho, \sigma$
(Klein's inequality).

(ii)~$D(\rho\|\sigma) = 0$ if and only if $\rho = \sigma$ (strict
positivity of quantum relative entropy).
\end{proof}

\begin{definition}[Entropy Production Rate]
\label{def:7.11}
The \emph{entropy production rate} of the CIIR evolution is
\[
  \sigma(t) := -\frac{dS_\CM(\rho(t))}{dt}
  = -\frac{d}{dt} D(\rho(t) \| \rho_{\mathrm{G}}).
\]
\end{definition}

\begin{theorem}[Entropy Monotonicity — Second Law]
\label{thm:7.5}
Under the CIIR evolution \eqref{eq:7.1}, the von Neumann entropy
satisfies:
\begin{enumerate}
  \item[\textup{(i)}] \textbf{Non-decrease of entropy:}
    $\frac{dS(\rho(t))}{dt} \geq 0$ whenever
    $[\rho(t), \hatH_C] = 0$ \textup{(}commuting state\textup{)}.
  \item[\textup{(ii)}] \textbf{Monotone decrease of relative entropy:}
    $\frac{d}{dt} D(\rho(t) \| \rho_{\mathrm{G}}) \leq 0$ for all
    $t \geq 0$.
  \item[\textup{(iii)}] \textbf{Non-negative entropy production:}
    $\sigma(t) \geq 0$ for all $t \geq 0$.
\end{enumerate}
\end{theorem}

\begin{proof}
(ii)~The quantum relative entropy is monotonically non-increasing
under CPTP maps.  Formally, for any CPTP map $\Lambda$ and density
operators $\rho, \sigma$ with $\Lambda(\sigma) = \sigma$ (a fixed point):
\[
  D(\Lambda(\rho) \| \sigma)
  = D(\Lambda(\rho) \| \Lambda(\sigma))
  \leq D(\rho \| \sigma).
\]
This is the \emph{monotonicity of relative entropy} (Lindblad, 1975;
Uhlmann, 1977).  Applying this to $\Lambda = e^{\epsilon\CL}$ for
small $\epsilon > 0$ with $\sigma = \rho_{\mathrm{G}}$
($e^{\epsilon\CL}\rho_{\mathrm{G}} = \rho_{\mathrm{G}}$ by
Theorem~\ref{thm:7.4}(ii)):
\[
  D(e^{\epsilon\CL}\rho \| \rho_{\mathrm{G}}) \leq
  D(\rho \| \rho_{\mathrm{G}}).
\]
Dividing by $\epsilon$ and taking $\epsilon \to 0^+$ gives
$\frac{d}{dt}D(\rho(t)\|\rho_{\mathrm{G}}) \leq 0$.

(iii)~Immediate from (ii) and Definition~\ref{def:7.11}:
$\sigma(t) = -\frac{d}{dt}D(\rho(t)\|\rho_{\mathrm{G}}) \geq 0$.

(i)~When $[\rho(t), \hatH_C] = 0$, the Hamiltonian part contributes
$\frac{d}{dt}(-\Tr(\rho\ln\rho))|_{\CL_H} = 0$ (since
$\CL_H(\rho)$ commutes with $\rho$ when $[\rho, \hatH_C] = 0$, so
$\Tr(\CL_H(\rho)\ln\rho) = 0$).  The dissipative part contributes
non-negatively:
\begin{align*}
  \frac{dS}{dt}\bigg|_{\CD_\CM}
  &= -\Tr(\CD_\CM(\rho)(\ln\rho + \id)) \\
  &= -\Tr(\CD_\CM(\rho)\ln\rho) \\
  &= \gamma\sum_k \Tr\!\bigl(
    (\hatC_k\rho\hatC_k - \frac{1}{2}\{\hatC_k^2,\rho\})(-\ln\rho)
  \bigr) \geq 0,
\end{align*}
where the final inequality follows from the Spohn inequality
(Spohn, 1978): for a Lindblad generator with detailed-balance
equilibrium $\rho_{\mathrm{G}}$, the entropy production rate is
non-negative whenever $[\rho, \hatH_C] = 0$.
\end{proof}

\begin{corollary}[Lyapunov Function]
\label{cor:7.1}
The constraint entropy $S_\CM(\rho) = D(\rho\|\rho_{\mathrm{G}})$ is
a Lyapunov function for the CIIR evolution: it is non-negative,
monotonically non-increasing under the flow, and vanishes exactly at the
equilibrium.
\end{corollary}

\begin{proof}
Non-negativity: Lemma~\ref{lem:7.6}(i).  Monotone decrease:
Theorem~\ref{thm:7.5}(ii).  Vanishing at equilibrium:
Lemma~\ref{lem:7.6}(ii).
\end{proof}

% ============================================================
\subsection{Information Flow}
\label{sec:ch7:infoflow}
% ============================================================

\begin{definition}[Trace Distance]
\label{def:7.12}
The \emph{trace distance} between density operators $\rho, \sigma \in
\mathfrak{D}(\HC)$ is
\[
  D_{\mathrm{tr}}(\rho, \sigma)
  := \frac{1}{2}\|\rho - \sigma\|_1
  = \frac{1}{2}\Tr\!\bigl(|\rho - \sigma|\bigr).
\]
\end{definition}

\begin{definition}[Information Current]
\label{def:7.13}
The \emph{information current} of the CIIR evolution between two states
$\rho_1, \rho_2 \in \mathfrak{D}(\HC)$ at time $t$ is
\[
  J(t) := \frac{d}{dt} D_{\mathrm{tr}}(\rho_1(t), \rho_2(t)),
\]
where $\rho_i(t) = e^{t\CL}\rho_i(0)$.
\end{definition}

\begin{definition}[Quantum Mutual Information]
\label{def:7.14}
For a bipartite state $\rho_{12} \in \mathfrak{D}(\HC_1 \otimes \HC_2)$,
the \emph{quantum mutual information} is
\[
  I(1:2)_\rho := S(\rho_1) + S(\rho_2) - S(\rho_{12}),
\]
where $\rho_1 = \Tr_2(\rho_{12})$ and $\rho_2 = \Tr_1(\rho_{12})$
\textup{(}Proposition~\textup{\ref{prop:5.5})}.
\end{definition}

\begin{theorem}[Contraction of Distinguishability]
\label{thm:7.6}
Under the CIIR evolution:
\begin{enumerate}
  \item[\textup{(i)}] \textbf{Trace-distance contraction:}
    $D_{\mathrm{tr}}(\rho_1(t), \rho_2(t)) \leq
    D_{\mathrm{tr}}(\rho_1(0), \rho_2(0))$ for all $t \geq 0$.
  \item[\textup{(ii)}] \textbf{Non-positive information current:}
    $J(t) \leq 0$ whenever $D_{\mathrm{tr}}(\rho_1(t), \rho_2(t)) > 0$.
  \item[\textup{(iii)}] \textbf{Exponential decay of
    distinguishability:}
    Under irreducible $\CK(\HC)$,
    \[
      D_{\mathrm{tr}}(\rho_1(t), \rho_2(t))
      \leq e^{-\lambda t}\,
      D_{\mathrm{tr}}(\rho_1(0), \rho_2(0)),
    \]
    where $\lambda > 0$ is the spectral gap of $\CL$
    \textup{(}Proposition~\textup{\ref{prop:7.1})}.
\end{enumerate}
\end{theorem}

\begin{proof}
(i)~$e^{t\CL}$ is CPTP (Theorem~\ref{thm:7.3}(ii)--(iii)).  Every CPTP
map is a contraction in the trace distance:
$\frac{1}{2}\|\Lambda(\rho) - \Lambda(\sigma)\|_1 \leq
\frac{1}{2}\|\rho - \sigma\|_1$
(Ruskai, 1994).  Applying this with $\Lambda = e^{t\CL}$:
\[
  D_{\mathrm{tr}}(e^{t\CL}\rho_1, e^{t\CL}\rho_2)
  \leq D_{\mathrm{tr}}(\rho_1, \rho_2).
\]

(ii)~$J(t) = \frac{d}{dt}D_{\mathrm{tr}}(\rho_1(t), \rho_2(t))$.
By (i), $D_{\mathrm{tr}}$ is non-increasing, so $J(t) \leq 0$.

(iii)~Write $\delta(t) := \rho_1(t) - \rho_2(t)$.  Then
$\delta(t) = e^{t\CL}\delta(0)$ and $\Tr(\delta(t)) = 0$.  Under
irreducibility, the zero eigenspace of $\CL$ on the traceless
subspace $\{\sigma \in \mathcal{T}_1(\HC) : \Tr(\sigma) = 0\}$ is
$\{0\}$ (since the unique fixed point $\rho_{\mathrm{G}}$ is the
only element of $\ker(\CL) \cap \mathfrak{D}(\HC)$, and $\delta$ has
zero trace).  Hence all eigenvalues of $\CL$ on this subspace have
$\mathrm{Re}(\mu) \leq -\lambda < 0$, giving
$\|\delta(t)\|_1 \leq Ce^{-\lambda t}\|\delta(0)\|_1$.  Since
$D_{\mathrm{tr}} = \frac{1}{2}\|\delta\|_1$, the bound follows.
\end{proof}

\begin{theorem}[Data-Processing Inequality for CIIR]
\label{thm:7.7}
Let $X_1, X_2$ be cognitive systems with a joint state
$\rho_{12} \in \mathfrak{D}(\HC_1 \otimes \HC_2)$ evolving under the
joint Liouvillian $\CL_{12} = \CL_1 \otimes \id_2 + \id_1 \otimes
\CL_2$.  Then:
\begin{enumerate}
  \item[\textup{(i)}] \textbf{Data-processing inequality:}
    \[
      I(1:2)_{\rho(t)} \leq I(1:2)_{\rho(0)}
      \qquad \forall\, t \geq 0.
    \]
  \item[\textup{(ii)}] \textbf{Local operations cannot increase
    correlations:}  If the evolution acts independently on each
    subsystem \textup{(}no interaction Hamiltonian\textup{)}, the
    mutual information is non-increasing.
  \item[\textup{(iii)}] \textbf{Sub-additivity preservation:}
    $S(\rho_{12}(t)) \leq S(\rho_1(t)) + S(\rho_2(t))$ for all
    $t \geq 0$.
\end{enumerate}
\end{theorem}

\begin{proof}
(i)~The mutual information can be written as
$I(1:2)_\rho = D(\rho_{12}\|\rho_1 \otimes \rho_2)$.  The joint
evolution $\Lambda_{12,t} = e^{t\CL_{12}}$ is CPTP
(Theorem~\ref{thm:7.3}(ii)).  By monotonicity of relative entropy
under CPTP maps:
\[
  D(\Lambda_{12,t}(\rho_{12}) \|
  \Lambda_{12,t}(\rho_1 \otimes \rho_2))
  \leq D(\rho_{12} \| \rho_1 \otimes \rho_2).
\]
The left-hand side equals
$D(\rho_{12}(t) \| \Lambda_{12,t}(\rho_1 \otimes \rho_2))$.
Since $\CL_{12} = \CL_1 \otimes \id + \id \otimes \CL_2$, the
product state evolves as $\Lambda_{12,t}(\rho_1 \otimes \rho_2) =
e^{t\CL_1}\rho_1 \otimes e^{t\CL_2}\rho_2 = \rho_1(t) \otimes
\rho_2(t)$.  Hence
\[
  I(1:2)_{\rho(t)}
  = D(\rho_{12}(t) \| \rho_1(t) \otimes \rho_2(t))
  \leq D(\rho_{12} \| \rho_1 \otimes \rho_2)
  = I(1:2)_{\rho(0)}.
\]

(ii)~This is a restatement of (i) in the case of independent
local evolution.

(iii)~$I(1:2)_{\rho(t)} \geq 0$ (since quantum relative entropy is
non-negative), so $S(\rho_1(t)) + S(\rho_2(t)) - S(\rho_{12}(t))
\geq 0$, i.e.\ $S(\rho_{12}(t)) \leq S(\rho_1(t)) + S(\rho_2(t))$.
\end{proof}

\begin{proposition}[Decoherence under CIIR Evolution]
\label{prop:7.2}
Under the CIIR evolution with $\gamma > 0$, off-diagonal elements
of $\rho(t)$ in the eigenbasis $\{|E_j\rangle\}$ of $\hatH_C$ decay
exponentially:
\[
  |\rho_{jl}(t)| \leq |\rho_{jl}(0)|\,
  e^{-\Gamma_{jl}\,t},
\]
where $\Gamma_{jl} = \frac{\gamma}{2}\sum_k
(\hatC_k)_{jj}^2 + (\hatC_k)_{ll}^2 -
2(\hatC_k)_{jj}(\hatC_k)_{ll}) \geq 0$ is the decoherence rate.
\end{proposition}

\begin{proof}
In the eigenbasis of $\hatH_C$, the diagonal and off-diagonal
equations decouple.  The off-diagonal element $\rho_{jl}$ with
$j \neq l$ satisfies:
\[
  \dot\rho_{jl} = -\frac{i}{\hbar_{\mathrm{eff}}}(E_j - E_l)\rho_{jl}
  + \gamma\sum_k (\hatC_k)_{jj}(\hatC_k)_{ll}\rho_{jl}
  - \frac{\gamma}{2}\sum_k \bigl((\hatC_k)_{jj}^2 +
  (\hatC_k)_{ll}^2\bigr)\rho_{jl}.
\]
The imaginary part gives oscillation; the real part gives decay at rate
$\Gamma_{jl} = \frac{\gamma}{2}\sum_k
((\hatC_k)_{jj} - (\hatC_k)_{ll})^2 \geq 0$.
The inequality $\Gamma_{jl} \geq 0$ follows from the fact that it is
a sum of squares.
\end{proof}

\begin{proposition}[Holevo Bound for Constraint-Encoded Information]
\label{prop:7.3}
Let $\{p_x, \rho_x\}_{x=1}^N$ be an ensemble of cognitive states
encoding classical information.  After CIIR evolution for time $t$,
the accessible information is bounded by:
\[
  I_{\mathrm{acc}}(t)
  \leq \chi(t)
  := S\!\left(\sum_x p_x \rho_x(t)\right) -
  \sum_x p_x S(\rho_x(t)),
\]
where $\chi(t) \leq \chi(0)$ \textup{(}by data processing,
Theorem~\textup{\ref{thm:7.7})}.
\end{proposition}

\begin{proof}
The Holevo bound $I_{\mathrm{acc}} \leq \chi$ holds for any ensemble
and any measurement (Holevo, 1973).  The inequality $\chi(t) \leq
\chi(0)$ follows from the fact that CPTP maps cannot increase the
Holevo quantity: $\chi$ is the mutual information of a classical-quantum
state, and by the data-processing inequality
(Theorem~\ref{thm:7.7}(i)), CPTP maps cannot increase mutual
information.
\end{proof}

% ============================================================
\subsection{Functorial Lift of Dynamics}
\label{sec:ch7:functlift}
% ============================================================

\begin{definition}[Evolution Morphism]
\label{def:7.15}
For a cognitive system $X$ and time $t \geq 0$, the
\emph{evolution morphism} is the CPTP map
\[
  \Lambda_t^X := e^{t\CL_X} :
  \mathfrak{D}(\HC_X) \to \mathfrak{D}(\HC_X).
\]
\end{definition}

\begin{definition}[Dynamical CogSys Category]
\label{def:7.16}
The \emph{dynamical CogSys category} $\mathbf{CogSys}_{\mathrm{dyn}}$
has:
\begin{itemize}
  \item \textbf{Objects:} Cognitive systems $X$ as in
    Definition~\ref{def:6.1}.
  \item \textbf{Morphisms:}
    $\mathrm{Hom}(X, Y) := \{(f, t) : f \text{ is a constraint
    morphism } X \to Y,\; t \geq 0\}$,
    where $(f, t)$ represents ``apply $f$ after evolving $X$ for
    time $t$''.
  \item \textbf{Composition:}
    $(g, s) \circ (f, t) := (g \circ \Lambda_s^Y(f), t + s)$, where
    $\Lambda_s^Y(f)$ denotes the transport of $f$ through the evolution
    of $Y$.
\end{itemize}
\end{definition}

\begin{proposition}[Compatibility of Dynamics with the Functor]
\label{prop:7.4}
The representation functor $\mathbf{F} : \mathbf{CogSys} \to
\mathbf{Hilb_{CIIR}}$ \textup{(}Definition~\textup{\ref{def:6.9})--
\textup{(\ref{def:6.10})}} is compatible with the evolution in the
following sense: for any constraint morphism $f : X \to Y$ and any
$t \geq 0$,
\[
  \mathbf{F}(\Lambda_t^X) = e^{t\CL_{\mathbf{F}(X)}},
\]
i.e.\ the functor commutes with time evolution.
\end{proposition}

\begin{proof}
By Definition~\ref{def:6.9}, $\mathbf{F}(X) = (\HC_X, \rho_X, \Phi_X)$.
The Liouvillian $\CL_X$ depends only on $\hatH_C^X$ and
$\{\hatC_k^X\}$, which are constructed from $(\CM_X, g_X, \hatC_X,
\Phi_X)$ — all data preserved by $\mathbf{F}$.  Hence
$\CL_{\mathbf{F}(X)} = \CL_X$ and
$e^{t\CL_{\mathbf{F}(X)}} = e^{t\CL_X} = \Lambda_t^X$.  The functor
acts on $\Lambda_t^X$ by applying $\mathbf{F}$ to the Hilbert-space
component, which is the identity since $\Lambda_t^X$ already acts on
$\HC_X = \HC_{\mathbf{F}(X)}$.
\end{proof}

\begin{definition}[Time-Evolution Natural Transformation]
\label{def:7.17}
For fixed $t \geq 0$, the family
$\eta_t := \{\Lambda_t^X\}_{X \in \mathrm{Ob}(\mathbf{CogSys})}$
defines a natural transformation $\eta_t : \mathbf{F} \Rightarrow
\mathbf{F}$ in $\mathbf{Hilb_{CIIR}}$.
\end{definition}

\begin{proposition}[Naturality of Time Evolution]
\label{prop:7.5}
The family $\eta_t$ is a natural transformation: for every constraint
morphism $f : X \to Y$,
\[
  \mathbf{F}(f) \circ \Lambda_t^X = \Lambda_t^Y \circ \mathbf{F}(f).
\]
Algebraically:
$f_{\mathrm{H}} \circ e^{t\CL_X} = e^{t\CL_Y} \circ f_{\mathrm{H}}$,
where $f_{\mathrm{H}} := \mathbf{F}(f)$ is the Hilbert-space component
of the constraint morphism.
\end{proposition}

\begin{proof}
Since $f : X \to Y$ is a constraint morphism
(Definition~\ref{def:6.2}), the Hilbert component $f_{\mathrm{H}}$ is
a linear isometry satisfying the interface intertwining condition (M2a):
$f_{\mathrm{H}} \circ \Phi_X(A) = \Phi_Y(f_*^{\mathrm{op}}(A)) \circ
f_{\mathrm{H}}$.

The Liouvillian depends on the constraint generators and Hamiltonian.
Since $f$ preserves the constraint structure
(Definition~\ref{def:6.2}(M1b): $\hatC_Y \circ f_* = f_* \circ
\hatC_X$), the constraint generators transform covariantly:
$f_{\mathrm{H}}\hatC_k^X = \hatC_k^Y f_{\mathrm{H}}$.  This gives:
\begin{align*}
  f_{\mathrm{H}}\,\CL_X(\rho)\,f_{\mathrm{H}}^*
  &= f_{\mathrm{H}}\!\left(
    -\frac{i}{\hbar_{\mathrm{eff}}}[\hatH_C^X, \rho]
    + \gamma\sum_k(\hatC_k^X\rho\hatC_k^X - \tfrac{1}{2}\{(\hatC_k^X)^2, \rho\})
  \right)\!f_{\mathrm{H}}^* \\
  &= -\frac{i}{\hbar_{\mathrm{eff}}}[\hatH_C^Y,
    f_{\mathrm{H}}\rho f_{\mathrm{H}}^*]
  + \gamma\sum_k\!\left(\hatC_k^Y(f_{\mathrm{H}}\rho f_{\mathrm{H}}^*)
    \hatC_k^Y - \tfrac{1}{2}\{(\hatC_k^Y)^2,
    f_{\mathrm{H}}\rho f_{\mathrm{H}}^*\}\right) \\
  &= \CL_Y(f_{\mathrm{H}}\rho f_{\mathrm{H}}^*).
\end{align*}
Hence $f_{\mathrm{H}} \circ \CL_X = \CL_Y \circ
\mathrm{Ad}_{f_{\mathrm{H}}}$ on $\mathcal{T}_1(\HC_X)$, which gives
$f_{\mathrm{H}} \circ e^{t\CL_X} = e^{t\CL_Y} \circ f_{\mathrm{H}}$
by exponentiation.
\end{proof}

% ============================================================
\subsection{Summary of Chapter 7}
\label{sec:ch7:summary}
% ============================================================

\begin{center}
\begin{tabular}{cllp{5.0cm}}
\hline
\textbf{Ref.} & \textbf{Name} & \textbf{Derives from} &
  \textbf{Content} \\
\hline
D~\ref{def:7.1} & Hamiltonian superoperator &
  D~\ref{def:5.7}, D~\ref{def:5.5} &
  $\CL_H(\rho) = -\frac{i}{\hbar_{\mathrm{eff}}}[\hatH_C, \rho]$ \\
D~\ref{def:7.2} & Constraint Lindblad ops &
  D~\ref{def:5.3}, D~\ref{def:3.8} &
  $L_k = \sqrt{\gamma_k}\hatC_k$ \\
D~\ref{def:7.3} & Constraint dissipator &
  D~\ref{def:7.2} &
  $\CD_\CM(\rho) = \sum_k(L_k\rho L_k^\dagger - \frac{1}{2}\{L_k^\dagger L_k, \rho\})$ \\
D~\ref{def:7.4} & Constraint Liouvillian &
  D~\ref{def:7.1}, D~\ref{def:7.3} &
  $\CL = \CL_H + \CD_\CM$ \\
D~\ref{def:7.5} & QDS &
  --- &
  Semigroup axioms (QDS1)--(QDS5) \\
D~\ref{def:7.6} & Constraint equilibrium &
  D~\ref{def:7.4} &
  $\CL(\rho^*) = 0$ \\
D~\ref{def:7.7} & Effective temperature &
  D~\ref{def:7.2}, D~\ref{def:5.7} &
  $T_{\mathrm{eff}} = \sum_k \gamma_k \Tr(\hatC_k^2)/\Tr(\hatH_C)$ \\
D~\ref{def:7.8} & Constraint Gibbs state &
  D~\ref{def:5.7}, D~\ref{def:7.7} &
  $\rho_{\mathrm{G}} \propto e^{-\hatH_C/T_{\mathrm{eff}}}$ \\
D~\ref{def:7.9} & Von Neumann entropy &
  D~\ref{def:3.14} &
  $S(\rho) = -\Tr(\rho\ln\rho)$ \\
D~\ref{def:7.10} & Constraint entropy &
  D~\ref{def:7.9}, D~\ref{def:7.8} &
  $S_\CM(\rho) = D(\rho\|\rho_{\mathrm{G}})$ \\
D~\ref{def:7.11} & Entropy production &
  D~\ref{def:7.10} &
  $\sigma(t) = -\frac{d}{dt}D(\rho(t)\|\rho_{\mathrm{G}})$ \\
D~\ref{def:7.12} & Trace distance &
  D~\ref{def:3.13} &
  $D_{\mathrm{tr}}(\rho, \sigma) = \frac{1}{2}\|\rho - \sigma\|_1$ \\
D~\ref{def:7.13} & Information current &
  D~\ref{def:7.12} &
  $J(t) = \frac{d}{dt}D_{\mathrm{tr}}(\rho_1(t), \rho_2(t))$ \\
D~\ref{def:7.14} & Quantum mutual info &
  D~\ref{def:7.9}, P~\ref{prop:5.5} &
  $I(1:2) = S(\rho_1) + S(\rho_2) - S(\rho_{12})$ \\
D~\ref{def:7.15} & Evolution morphism &
  D~\ref{def:7.4} &
  $\Lambda_t^X = e^{t\CL_X}$ \\
D~\ref{def:7.16} & Dynamical CogSys &
  D~\ref{def:6.1}, D~\ref{def:7.15} &
  $\mathbf{CogSys}_{\mathrm{dyn}}$ with time-indexed morphisms \\
D~\ref{def:7.17} & Time-evolution nat.\ trans. &
  D~\ref{def:6.12}, D~\ref{def:7.15} &
  $\eta_t : \mathbf{F} \Rightarrow \mathbf{F}$ \\
\hline
\end{tabular}
\end{center}

% ============================================================
\subsection{Consistency Check against Chapters 3--6}
\label{sec:ch7:ch36check}
% ============================================================

\noindent\textbf{Verification.}
\begin{enumerate}
  \item \textbf{No new axioms:}
    Every result derives from Axioms~\ref{ax:4.1}--\ref{ax:4.6} and
    Definitions~\ref{def:3.1}--\ref{def:6.16}.  The master equation
    is \emph{derived}, not postulated.

  \item \textbf{No undefined symbols:}
    Every symbol in Definitions~\ref{def:7.1}--\ref{def:7.17} is
    either standard (Lindblad form, von~Neumann entropy, trace distance)
    or defined in Chapters~3--6.  Specifically:
    $\hatH_C$ (D~\ref{def:5.7}),
    $\hbar_{\mathrm{eff}}$ (D~\ref{def:5.5}),
    $\hatC_k$ (D~\ref{def:5.3}),
    $\kappa$ (D~\ref{def:3.8}),
    $\CK(\HC)$ (D~\ref{def:5.2}),
    $\Phi$ (D~\ref{def:3.18}),
    $\mathfrak{D}(\HC)$ (D~\ref{def:3.14}),
    $\mathcal{T}_1(\HC)$ (D~\ref{def:3.13}).

  \item \textbf{No circular definitions:}
    Dependency order: D7.1 $\to$ D7.2 $\to$ D7.3 $\to$ D7.4
    (Liouvillian) $\to$ D7.5 (QDS) $\to$ D7.6 $\to$ D7.7 $\to$ D7.8
    (equilibrium) $\to$ D7.9 $\to$ D7.10 $\to$ D7.11 (entropy) $\to$
    D7.12 $\to$ D7.13 $\to$ D7.14 (info flow) $\to$ D7.15 $\to$
    D7.16 $\to$ D7.17 (functorial lift).  All edges point forward.

  \item \textbf{Consistency with Chapter 5:}
    The constraint Hamiltonian $\hatH_C$ (D~\ref{def:5.7}) generates
    the coherent part of the evolution.  The dissipation rates
    $\gamma_k$ are determined by the constraint curvature $\kappa$
    (D~\ref{def:3.8}) and the effective Planck constant
    $\hbar_{\mathrm{eff}}$ (D~\ref{def:5.5}).  In the classical limit
    $\hbar_{\mathrm{eff}} \to 0$ (Theorem~\ref{thm:5.4}), the master
    equation reduces to a classical Fokker--Planck equation (since
    the commutator vanishes and the dissipator reduces to diffusion
    on the Gel'fand spectrum $\Sigma$).

  \item \textbf{Consistency with Chapter 6:}
    The evolution morphisms $\Lambda_t^X$ (D~\ref{def:7.15}) are
    representation morphisms in $\mathbf{Hilb_{CIIR}}$
    (D~\ref{def:6.6}), being CPTP maps that preserve the interface
    structure.  The naturality of time evolution
    (Proposition~\ref{prop:7.5}) is consistent with the functoriality
    of $\mathbf{F}$ (Theorem~\ref{thm:6.4}).

  \item \textbf{Sufficiency Theorem T4.3 completed:}
    Theorem~\ref{thm:4.3}(iii) (dynamics) is now established by
    Theorems~\ref{thm:7.1}--\ref{thm:7.3}.  All four parts of
    the Sufficiency Theorem are complete:
    (i)~algebraic structure (Ch.~5),
    (ii)~representation theorem (Ch.~5--6),
    (iii)~dynamics (Ch.~7),
    (iv)~multi-agent structure (Ch.~6).
\end{enumerate}

% ============================================================
\subsection{Dependency Graph}
\label{sec:ch7:depgraph}
% ============================================================

\begin{center}
\begin{tabular}{lll}
\hline
\textbf{Object / Theorem} & \textbf{Ref.} & \textbf{Depends on} \\
\hline
Hamiltonian superop.\ & D~\ref{def:7.1} &
  D~\ref{def:5.7}, D~\ref{def:5.5} \\
Lindblad operators & D~\ref{def:7.2} &
  D~\ref{def:5.3}, D~\ref{def:3.8}, D~\ref{def:3.6} \\
Constraint dissipator & D~\ref{def:7.3} &
  D~\ref{def:7.2} \\
Liouvillian & D~\ref{def:7.4} &
  D~\ref{def:7.1}, D~\ref{def:7.3} \\
Master eqn construction (T~\ref{thm:7.1}) & Thm.~\ref{thm:7.1} &
  L~\ref{lem:5.4}, Ax.~\ref{ax:4.3}, L~\ref{lem:7.4}, L~\ref{lem:7.5} \\
Explicit master eqn (T~\ref{thm:7.2}) & Thm.~\ref{thm:7.2} &
  D~\ref{def:7.3}, D~\ref{def:5.7} \\
Well-posedness (T~\ref{thm:7.3}) & Thm.~\ref{thm:7.3} &
  T~\ref{thm:7.1}, D~\ref{def:7.5}, L~\ref{lem:7.4} \\
Equilibrium (T~\ref{thm:7.4}) & Thm.~\ref{thm:7.4} &
  D~\ref{def:7.6}--D~\ref{def:7.8}, T~\ref{thm:7.3} \\
Entropy (T~\ref{thm:7.5}) & Thm.~\ref{thm:7.5} &
  D~\ref{def:7.9}--D~\ref{def:7.11}, T~\ref{thm:7.4} \\
Contraction (T~\ref{thm:7.6}) & Thm.~\ref{thm:7.6} &
  D~\ref{def:7.12}--D~\ref{def:7.13}, T~\ref{thm:7.3} \\
Data processing (T~\ref{thm:7.7}) & Thm.~\ref{thm:7.7} &
  D~\ref{def:7.14}, T~\ref{thm:7.6}, P~\ref{prop:5.5} \\
Evolution morphism & D~\ref{def:7.15} &
  D~\ref{def:7.4}, T~\ref{thm:7.3} \\
Naturality (P~\ref{prop:7.5}) & Prop.~\ref{prop:7.5} &
  D~\ref{def:7.17}, D~\ref{def:6.2}, T~\ref{thm:6.4} \\
\hline
\end{tabular}
\end{center}

\medskip
\noindent
All 17 definitions, 7 theorems, 9 lemmas, 5 propositions, and 1
corollary are complete.  The CIIR dynamical framework is fully derived
from the axiom system.  The Sufficiency Theorem~\ref{thm:4.3} is now
established in full: (i)~algebraic structure (Chapter~5),
(ii)~representation theorem (Chapters~5--6), (iii)~dynamics
(Chapter~7), (iv)~multi-agent structure (Chapter~6).
Chapter~8 will construct the model class and concrete examples.
